\section{Vulnerability Analyses}
\label{sec:vulnerability}
\subsection{Reverse Engineering of Magnetic Tactile Sensors}
\label{sec:reverse engineering}
To address the aforementioned questions, it is critical to understand the specifications of COTS magnetic tactile sensors. 
We propose an offline, model-specific \emph{reverse engineering approach}. This process can be performed on a separately obtained sensor of the same model and requires no access to the robotic system.
Specifically, we construct a three-dimensional testing space on the contact interface of the magnetic tactile sensor, and utilize small passive magnets to incur the interference.
This procedure reveals how the tactile sensor would generate force readings given various magnetic field settings.

As shown in Fig.~\ref{fig:vulnerability1}, we used the PaXini palm tactile sensor (PX6AX-GEN3-MC-M2020-Elite) for reverse engineering because of its resilience to magnetic field interference.
With the geometric center of the sensor's contact interface as the origin point, we defined a three-dimensional testing space.
The plane of this testing space was partitioned using grid paper with 2 mm × 2 mm squares, and its coverage area is 12 mm × 12 mm, i.e., ranging from -6 mm to 6 mm along the X and Y axes, with a planar step size of 2 mm.
The vertical space was partitioned using a vernier caliper (Nscing Es 889-101~\cite{ebuy7_nscing_889101}, precision 0.01 mm).
The vertical height ranges from 0 mm to 2.0 mm along the Z axis, with a vertical step size of 0.25 mm.
With this procedure, we set a 12 mm × 12 mm × 2 mm three-dimensional testing space above the sensor's contact interface.
These settings are depicted in Fig.~\ref{fig:vulnerability1}.
For the magnetic interferences, we used two specifications of cylindrical magnets.
Specifically, we used N35-grade passive magnets with a thickness of 0.5 mm and diameters of 2.5 mm and 3 mm.
Next, we moved the magnet to different positions atop the tactile sensor and recorded the force readings, i.e., ($F_x, F_y, F_z$).
This process yields a dataset of 882 discrete force-feedback measurements.
Based on this dataset, we constructed spatial heatmaps for the three-dimensional forces. 
As shown in Fig.~\ref{fig:vulnerability2}, these heatmaps illustrate the distribution of the tactile sensor's response to external magnetic stimuli. 

\begin{figure}[t] 
  \centering
  \includegraphics[width=\linewidth]{figures/vulnerability1.pdf}
  \Description{Logical diagram illustrating the reverse engineering process of PhantomTac, including the magnetic tactile sensor, measurement setup, and the procedure used to characterize sensor responses for constructing the attack.}
  \caption{Setup for the reverse engineering process.}
  \label{fig:vulnerability1}
\end{figure}

\subsection{Design of Attack Modalities}
\begin{figure*}[htbp] 
  \centering
  \includegraphics[width=0.95\textwidth]{figures/vulnerability2.png} 
  \Description{Three-dimensional force feedback visualization of a magnetic tactile sensor, showing the sensor response under magnetic interference from a 3 mm diameter magnetic source.}
  \caption{Three-dimensional force feedback of magnetic tactile sensor. The diameter of the magnetic interference source is 3 mm.}
  \label{fig:vulnerability2}
\end{figure*}

After obtaining the response characteristics of the magnetic tactile sensor to external magnetic stimulation, we further conducted the attack design of the magnetic tag. 
To achieve an effective attack on the magnetic tactile sensor, the tag design needs to simultaneously satisfy the following conditions:

\noindent\textbf{(1) Balancing stealth and the magnetic field strength.} 
The tag must maintain a minimal form factor (size and thickness) to avoid visual detection or mechanical interference during deployment. 
However, the volume must remain sufficient to generate a magnetic flux density capable of inducing a measurable tactile response.

\noindent\textbf{(2) Decoupled force injection.} The design must allow for the independent manipulation of $F_x, F_y,$ and $F_z$.
The robotic control system interprets shear forces ($F_x$ and $F_y$) and normal forces ($F_z$) differently.
For example, a robotic gripper can use the shear force variations for detecting slip.
Therefore, the attacker may wish to inject significant tangential anomalies while suppressing $F_z$.

To achieve this, we identify high-sensitivity regions for specific force components. By strategically positioning a miniature magnet at a coordinate where the sensor exhibits a strong response to the target component and minimal response to others, we maximize the fidelity of the injected force signal.
While heatmaps provide a coarse distribution, we perform spatial interpolation modeling in Sec.~\ref{sec:design} to derive continuous response functions, enabling the precise localization of optimal injection points.

As we will show in Sec.~\ref{sec:design}, with this design principle, we can inject a normal force of approximately 4.2 N into the magnetic tactile sensor with only one thin passive magnet (with diameter of 3 mm). This magnitude of false force feedback is already sufficient to cause gripping force regulation anomalies in real-world robotic systems, such as excessively loose or excessively tight grasping. Meanwhile, the impact of this magnet on the force components in the other two directions does not exceed 0.7 N.

\subsection{Improving the Attack Stealthiness}
To achieve the \emph{deploy-and-forget} capability, it is essential to ensure stealthiness. 
In {\name}, we selected miniature magnets with a thickness of 0.5 mm as the attack source. 
Miniature magnets of this size can be concealed beneath camouflage labels without exhibiting obvious visual anomalies.

\begin{table}[t]
    \centering
    \caption{Common Sticker Thicknesses.}
    \label{tab:label_thickness}
    \resizebox{\columnwidth}{!}{%
    \begin{tabular}{llc}
        \toprule
        \textbf{Label Type} & \textbf{Common Applications} & \textbf{Thickness (mm)} \\ 
        \midrule
        Thermal Paper~\cite{avery_direct_thermal_200gp} & Express shipping labels & $0.142 \pm 10\%$ \\
        Coated Paper~\cite{avery_high_gloss_paper} & Product barcodes & $0.183 \pm 10\%$ \\
        Synthetic Paper~\cite{avery_synthetic_paper_bw0153} & Waterproof QR code labels & $0.145 \pm 10\%$ \\
        PET Label~\cite{avery_pet_sy7201} & Electronic product back labels & $0.115 \pm 10\%$ \\
        \bottomrule
    \end{tabular}%
    }
\end{table}

As shown in Table~\ref{tab:label_thickness}, common logistics and product labels typically range from $0.1$ to $0.2$~mm in thickness. 
Because the non-ferromagnetic substrates (e.g., paper, adhesive, PET) have a relative magnetic permeability near unity ($\mu_r \approx 1$), 
this thin camouflage does not impede the efficacy of the magnetic tag underneath.

We performed an experiment to verify this. 
Specifically, we performed a vertical scan using the 2.5 mm diameter magnet --- our smallest one. 
By constraining the planar coordinates $(x, y)$ and treating the $Z$-axis elevation as the variable, we profiled the attenuation profiles of the induced $F_x, F_y$, and $F_z$ components. 
Based on the spatial heatmaps (Fig.~\ref{fig:vulnerability2}), we selected $(0, 0, z)$, $(0, 2, z)$, and $(2, 0, z)$ as the data extraction trajectories to characterize the force-distance relationship.

\begin{figure}[t] 
  \centering
  \includegraphics[width=0.97\linewidth]{figures/vulnerability3.pdf}
  \Description{Experimental evaluation of the penetration capability of a 2.5 mm diameter magnet, showing the effect of magnetic interference through different distances or material layers on the tactile sensing response.}
  \caption{The penetration capability of 2.5 mm magnet.}
  \label{fig:vulnerability3}
\end{figure}

\begin{figure}[t] 
  \centering
  \includegraphics[width=1\linewidth]{figures/vulnerability4.pdf}
  \Description{Examples of disguised appearances of PhantomTac tags. The subfigures show the attack device embedded or concealed as a QR code, a blank paper label, a shipping label, a cartoon sticker, and a product sticker.}
  \caption{Disguised appearance of {\name}. (a)---(e) present QR code, blank paper, shipping label, cartoon sticker, and product sticker, respectively.}
  \label{fig:vulnerability4}
\end{figure}

As shown in Fig.~\ref{fig:vulnerability3}, the magnet maintains a stable signal penetration at a vertical offset of $1$ mm,
thus significantly exceeding the thickness of standard labels. 
Note that higher-grade magnets (e.g., N52) would further extend this effective range. 
As demonstrated in Fig.~\ref{fig:vulnerability4}, the resulting {\name} tags are indistinguishable from ordinary stickers, allowing them to remain visually concealed before robotic contact. 
This stealthiness distinguishes {\name} from conventional sabotage (e.g., applying a slippery substance or directly damaging the object). 
Specifically, although conventional sabotage can lead to direct damage, this approach requires unrestricted physical access and can be easily noticed by the victim. 
In contrast, {\name} can instead remain concealed until the tagged object contacts the tactile sensor. 
Therefore, even if the tag is discovered after an incident, object damage or workflow disruption has already occurred.
